\section{Kiểm tra T38: dịch vụ, miễn trừ và biên nhìn thấy}
\label{muc:t38}
\label{muc:biendukien-t38}

\subsection{Thiết kế được cố định trước tập kiểm tra}

T38 chuyển từ tập ô phát triển sang ngày 14--18/11, bốn mẫu mỗi ngày, ba khung
07:00--09:00, 08:00--10:00 và 17:00--19:00. Nó hỏi hai điều: cổng hành vi còn
giữ số khách gần các cổng khác khi phục hồi hay không, và một khoảng dự phòng
ước lượng trước có bao phần nhìn thấy trong các ô còn lại hay không. Mỗi ô
vẫn có 108 yêu cầu, 8 xe và mạng nhỏ 96 nút. Biến thiên thời gian đi lấy từ
chuỗi thực; không có kịch bản khuếch đại tổng hợp trong đợt này.

Kế hoạch T38 được ghi trong Git ở commit \texttt{323a8d9}, trước lần hiệu chỉnh
đầu tiên ngày 01/10/2026. Tệp khoảng dự phòng được ghi ở commit
\texttt{19517b0}, trước lần kiểm tra đầu tiên. Những mốc này xác nhận thứ tự
cố định kế hoạch và ngưỡng, không biến thiết kế thành xác nhận độc lập:
các trần, biên chấp nhận và cách chia được chọn sau khi đã xem các đợt phát triển.
Tập ô ngày 14--18/11 cũng từng được dùng ở tầng 2 với phiên bản thế giới khác.
Trong khóa luận, T38 được gọi là \emph{thăm dò có kế hoạch cố định trước tập kiểm tra}.

Kế hoạch dự kiến 60 ô. Mười một ô không dựng được vì vượt trần số nút của
công cụ chuẩn hóa: bốn ô tối ngày 15/11, ba ô sáng 07:00 ngày 17/11 và bốn ô
sáng 07:00 ngày 18/11. Loại ô xảy ra ở bước dựng đầu vào, không theo kết quả
dịch vụ. Dù vậy, nó làm thay đổi thành phần ngày và khung giờ; kết quả không
đại diện cho toàn bộ 60 ô dự kiến.

Bốn mẫu của mỗi ngày được chia theo nhóm $(\text{ngày},\text{mẫu})$, với hạt
giống chia 20261001. Toàn bộ khung của cùng nhóm nằm cùng một phía. Có 15 ô
hiệu chỉnh thuộc sáu nhóm và 34 ô kiểm tra thuộc 14 nhóm. Số ô sáng 07:00,
sáng 08:00 và tối trong hiệu chỉnh là 4, 6, 5; trong kiểm tra là 9, 14, 11.
Một ô hiệu chỉnh đã được đọc trong phép chạy thử, nên không được gọi toàn
bộ tập hiệu chỉnh là dữ liệu chưa từng được xem.

\subsection{Giao nhau của chuyến nguồn và đơn vị phân tích}


Tách nhóm không đủ để khẳng định các chuyến ở hai phía rời nhau. Các tệp mẫu
nhu cầu khác nhau vẫn có thể chứa cùng chuyến taxi gốc. Lần biên tập này nối
108 bản ghi được chọn của mỗi ô về 20 tệp CSV gốc, dùng cặp ngày và
\texttt{request\_id} trong CSV để nhận diện chuyến, rồi kiểm thời điểm yêu cầu
và hai điểm đầu--cuối. Kết quả có \textbf{75 chuyến nguồn chung giữa hiệu chỉnh
và kiểm tra}; các trường gốc trùng nhau. Số chuyến riêng biệt là 1.574 ở hiệu
chỉnh và 3.505 ở kiểm tra. Số thứ tự dòng chỉ có nghĩa trong một tệp mẫu,
nên không được dùng riêng số đó làm định danh chuyến.

Giao nhau này không làm phép cộng số khách hoặc kiểm sổ cái sai. Nó làm yếu
cách diễn giải phần kiểm tra như dữ liệu độc lập với hiệu chỉnh. Các ô cũng
chia sẻ ngày, hệ số thời gian đi và một phần nhu cầu. Vì thế đơn vị báo cáo
là ô, đơn vị nhóm khi chia là ngày--mẫu, còn đơn vị khách trong các tổng không
phải các quan sát độc lập cho một kiểm định về quần thể.

Các khoảng 95\% được báo trong T38 dùng 10.000 lần lấy lại mẫu có hoàn lại
theo \emph{ô ghép cặp}, hạt giống 20261002. Cùng một chỉ số ô được lấy cho hai
nhánh trong mỗi lần. Cách này giữ cặp so sánh nhưng chưa xử lý phụ thuộc giữa
các ô cùng ngày hoặc chuyến nguồn. Khóa luận giữ chúng để đọc độ phân tán theo
quy tắc phân tích đã ghi; không coi đó là khoảng suy luận đã hiệu chỉnh theo
cụm hay bằng chứng tương đương trong quần thể. Không thêm kiểm định mới lên
dữ liệu chưa đáp ứng giả thiết độc lập.

\subsection{Số phục vụ gần nhau, nghĩa vụ khác nhau}

Tập kiểm tra có 34 ô, tức $34\times108=3.672$ yêu cầu mỗi nhánh. Mỗi ô chạy năm
nhánh: $U$, $C$-30, $M^{+}$-30, $M^{+}$-60 và $V$-30. Bảng~\ref{bang:t38-phucvu}
giữ tổng phục vụ và hai loại gánh khác nhau. Các tỉ lệ ở hai dòng giữa là
trung bình không trọng số của tỉ lệ từng ô; không phải gộp theo khách.

\begin{table}[htbp]
\centering
\caption{Kết quả T38 trên 34 ô kiểm tra. ``Vượt lời hứa riêng'' dùng đại lượng
của chính cổng; tỉ lệ quyết định miễn trừ có mẫu số khác tỉ lệ khách.}
\label{bang:t38-phucvu}
\small
\begin{tabularx}{\linewidth}{@{}Xrrrr@{}}
\toprule
\textbf{Đại lượng} & \textbf{$C$-30} & \textbf{$M^{+}$-30} & \textbf{$M^{+}$-60} & \textbf{$V$-30}\\
\midrule
Yêu cầu hoàn tất & 2.338 & 2.338 & 2.347 & 2.336\\
Tỉ lệ miễn trừ (\%) & 0,0 & 41,3 & 22,5 & 53,8\\
Khách vượt lời hứa riêng (\%) & 0,0 & 14,5 & 4,2 & 21,6\\
Ô có khách vượt lời hứa riêng & 0/34 & 28/34 & 22/34 & 34/34\\
\bottomrule
\end{tabularx}
\end{table}

Tổng của $C$-30 bằng $M^{+}$-30 và hơn $V$-30 hai khách. Chênh theo ô có thể
theo cả hai chiều, như Hình~\ref{hinh:t38-cap}. Trung bình $C-M^{+}$ là 0,000
khách mỗi ô, khoảng lấy lại mẫu $[-0{,}147;\,0{,}118]$; $C-V$ là 0,059,
khoảng $[-0{,}118;\,0{,}206]$. Chênh tuyệt đối lớn nhất ở một ô là hai khách.
Theo quy tắc đã đặt của đợt này, các khoảng nằm trong biên $\pm2$ khách mỗi ô.
Kết quả hỗ trợ sự gần nhau trong tập ô, không hỗ trợ ``$C$ phục vụ nhiều hơn''.

\begin{figure}[htbp]
\centering
\includegraphics[width=\linewidth]{hinh/t38_paired_final.pdf}
\caption{Chênh số hoàn tất ghép cặp trên 34 ô T38. Mỗi dấu biểu diễn một ô;
trùng vị trí là nhiều ô có cùng chênh. Không lấy số lần chạy chính sách làm
số mẫu độc lập. Chi phí so với $U$ phân tán hơn chênh giữa các cổng.}
\label{hinh:t38-cap}
\end{figure}

So với $U$, chi phí vẫn tồn tại. $U$ hoàn tất 2.434 yêu cầu, hơn $C$-30 96.
Chênh trên tổng yêu cầu là $96/3.672\times100=2{,}61$ \emph{điểm phần trăm};
nếu lấy số khách $U$ làm mẫu số thì giảm tương đối 3,94\%. Không được gọi
cả hai là cùng một ``2,6\%''. Chênh trung bình là 2,824 khách mỗi ô, khoảng
lấy lại mẫu $[1{,}794;\,3{,}824]$, nằm dưới trần chi phí dịch vụ năm khách mỗi ô của kế hoạch.
Trần này là lựa chọn thăm dò sau phát triển, không phải mức chi phí đã được
người vận hành hay hành khách chấp thuận.

Chi phí không đồng đều: $U-C$ trải từ $-4$ đến $+9$ khách, và 21/34 ô có độ
lớn chênh quá hai. Tổng thiếu của $C$ so với $U$ ở ba khung là 5, 39 và 52
khách. $M^{+}$-60 phục vụ hơn $C$-30 chín khách; đây là mô tả ngoài phép so
tương đương ở trần 30 giây, không được dùng để thay đổi giả thuyết sau khi
đọc kết quả.

Khác biệt lớn hơn nằm ở gánh phục hồi. Không quyết định nào của $C$ có miễn
trừ trên cả 49 ô hiệu chỉnh và kiểm tra; giá trị sổ giờ trả lớn nhất là
29,239 giây. Các cổng kia thường xuyên phải miễn trừ và vẫn có khách vượt
đại lượng đã hứa. Tuy nhiên, ``0\%'' của $C$ là kiểm tuân thủ một nghĩa vụ
hành vi hẹp, không phải tỉ lệ đến muộn bằng không. $C$-30 và $M^{+}$-30 có
cùng toàn bộ chuỗi hành động ở 30/34 ô; $C$-30 và $V$-30 ở 27/34. Không thể
suy ra cùng số phục vụ thì cùng khách được phục vụ hoặc cùng mọi giờ đến.

Nếu gộp theo khách, tỉ lệ vượt lời hứa riêng của $M^{+}$-30 là
$289/2.338=12{,}4\%$ và của $M^{+}$-60 là $79/2.347=3{,}4\%$,
khác với 14,5\% và 4,2\% trong bảng. Việc giữ mẫu số và cách gộp giúp tránh
biến một sự khác nhau về trọng số thành mâu thuẫn kết quả.

Chỉ số chung E2 có trung bình theo ô 5,78\% với $C$-30 và 5,70\% với
$M^{+}$-30 hoặc $V$-30. Chênh là 0,000857 theo thang tỉ lệ, tức 0,0857
điểm phần trăm; khoảng lấy lại mẫu là $[-0{,}0082;\,0{,}2218]$ điểm.
Nó nằm trong biên 0,02 theo thang tỉ lệ của kế hoạch, tức hai điểm phần trăm.
Đây là kiểm trên một thước đo chung, không là xác nhận về hài lòng của người dùng.

\subsection{Khoảng dự phòng từ phần thay đổi khi giữ tuyến}

Đối với mỗi ô hiệu chỉnh $j$, điểm số $S_j$ là tổng biến động giữ tuyến lớn
nhất theo khách đã được hứa, trên tọa độ giờ trả. Nếu một khách đã được hứa
chưa hoàn tất thì quy tắc gán điểm vô hạn; không có trường hợp đó trong 185
lần chạy hợp lệ. Lấy cực đại theo khách làm điểm của \emph{cả ô}, thay vì
hiệu chỉnh một tỉ lệ theo từng lần sửa ETA.

Quy tắc phân vị lấy cảm hứng từ dự báo conformal chia mẫu
\cite{angelopoulos2023} đặt $\alpha=0{,}1$, với $n=15$ ô hiệu chỉnh:
\begin{equation}
 k=\left\lceil(n+1)(1-\alpha)\right\rceil=15,\qquad
 \gamma=S_{(15)}=222{,}259\ \text{s}.
\end{equation}
Ở đây $S_{(15)}$ là điểm thứ 15 khi xếp tăng dần, tức điểm lớn nhất. Nó thuộc
ô ngày 14/11, mẫu 2, khung 07:00; điểm cao thứ hai là 165,116 giây. Một ô
quyết định toàn bộ khoảng dự phòng nên ước lượng khá nhạy với thành phần hiệu chỉnh.

Khi phần giữ tuyến của mọi khách trong ô không quá $\gamma$, trần hành vi
$\beta=30$ giây và bất đẳng thức ở Mục~\ref{muc:biendukien} cho
$V_r\le\beta+\gamma=252{,}259$ giây đối với chuỗi lời hứa giờ trả được đo.
Đây là biên có điều kiện cho \emph{biến động lời hứa}, không phải giới hạn
trễ đến nơi so với giờ hẹn đầu. Biên dùng tổng độ lớn nên còn tính cả cập
nhật sớm hơn; nó không thể thay một thỏa thuận mức dịch vụ
(Service Level Agreement, SLA) về giờ trả thực tế.

Trong dự báo conformal, bảo đảm độ phủ biên yêu cầu tính trao đổi được của
các điểm hiệu chỉnh và điểm mới. T38 chưa chứng minh giả thiết này; các ô
khác khung và chia sẻ ngày, chuyến nguồn. Vì thế công thức phân vị vẫn tính
được nhưng khóa luận không nhận bảo đảm 90\% cho quần thể. Thứ tự niêm phong
chỉ chứng minh không sửa $\gamma$ sau khi xem tập kiểm tra.

\subsection{Độ phủ quan sát và một trường hợp ngoài biên}

Hình~\ref{hinh:t38-diem} trình bày toàn bộ điểm của 15 ô hiệu chỉnh và 34 ô
kiểm tra. Hai điểm kiểm tra vượt $\gamma$: ngày 14/11, mẫu 3, khung 07:00
có 252,492 giây; ngày 15/11, mẫu 2, cùng khung có 226,400 giây. Tần suất
điểm không vượt ngưỡng là $32/34=94{,}1\%$. Riêng 07:00--09:00 là
$7/9=77{,}8\%$; các khung khác là 14/14 và 11/11. Độ phủ chung cao hơn
không loại được sự thiếu phủ tại khung sáng sớm.

\begin{figure}[htbp]
\centering
\includegraphics[width=\linewidth]{hinh/t38_scores_final.pdf}
\caption{Điểm số theo ô trong T38, tách hiệu chỉnh và kiểm tra. Đường ngưỡng
$\gamma=222{,}259$ giây được cố định trước tập kiểm tra. Điểm số là phần giữ
tuyến lớn nhất theo khách trong một ô, khác với biến động nhìn thấy $V$.}
\label{hinh:t38-diem}
\end{figure}

Vượt $\gamma$ không tự động có nghĩa $V$ vượt $\beta+\gamma$. Trong hai ô
vừa nêu, giá trị hành vi lớn nhất bằng không. Ô 226,400 giây vẫn nằm dưới
biên nhìn thấy 252,259 giây; ô 252,492 giây vượt biên đúng 0,233 giây.
Do đó có 32/34 ô được phủ theo \emph{điểm giữ tuyến}, nhưng 33/34 ô không có
khách vượt \emph{biên nhìn thấy}. Ở 32 ô được phủ theo điểm, $V$ lớn nhất là
216,005 giây và không có ngoại lệ đối với bất đẳng thức. Hai mẫu số giống
nhau không làm hai sự kiện này trở thành cùng một chỉ số.

Theo kế hoạch, tiêu chí không quá ba ô vượt điểm và biên không quá 300 giây
đều đạt. Tuy nhiên, 300 giây là ngưỡng hữu dụng chọn sau phát triển; biên
252,259 giây không đạt mốc tham chiếu 120 giây. Một khoảng rộng có thể dễ
phủ hơn nhưng ít giá trị hứa hẹn hơn. Khóa luận giữ cả hai mốc để người đọc
đánh giá, thay vì chỉ chọn mốc cho kết quả đạt.

Cận trên Clopper--Pearson một phía 95\% đã báo cho 2/34 là 17,4\%, không
nằm dưới 10\%. Cận này còn dựa trên mô hình nhị thức với các lần thử độc
lập; phụ thuộc giữa các ô khiến không thể dùng nó như cận đã được bảo đảm
cho thiết kế T38. Nó được giữ như giá trị danh nghĩa của phép tính ban đầu,
không là kết luận về tỉ lệ vi phạm trong quần thể.

\subsection{Độ tin cậy thực thi và điều được phép kết luận}

Tổng 185 lần chạy được dùng gồm 15 hiệu chỉnh và 170 kiểm tra. Hai lần chạy
lỗi hạ tầng được chạy lại, một do máy ngủ và một do quá hạn phản hồi; thư mục lỗi
được giữ. Phân tích chỉ dùng lần chạy lại hoàn tất theo quy tắc đã ghi.
Mỗi nhánh có cùng băm cấu hình trên mọi ô; các nhánh dùng cùng băm kiểm kê kho mã
(inventory) và cùng chuỗi hệ số đúng mốc 900 giây. Băm kiểm kê bao gồm commit Git,
đường dẫn, kích thước và nội dung các tệp được Git theo dõi hoặc chưa theo dõi nhưng
không bị bỏ qua. Bản kiểm độc lập đọc quyết định và sự
kiện xuống xe từ biên bản, rồi tái tính lịch sử ba mốc trong điểm kiểm tra
cuối. Đối chiếu $|x|$, $|c|$, $|z|$ ở 382.006 lần sửa giờ trả không có sai
khác trong phạm vi đó. Lần biên tập còn giải mã và tái kiểm băm các điểm
kiểm tra cuối, không chỉ chép lại bảng báo cáo.

Số phục vụ được đọc từ lịch sử vòng đời ở điểm kiểm tra cuối. Tệp tổng kết
không cung cấp một trường số phục vụ độc lập thứ hai; không thể gọi phép
đối chiếu đó là hai phép đo độc lập về dịch vụ. Kiểm băm bảo vệ tính toàn
vẹn của khung ghi, không chứng minh mô hình đã phản ánh giao thông thực.

Đợt này cho phép ghi nhận đón muộn tới 60 giây để tiếp tục mô phỏng. Có 43
lần đón muộn ở $C$-30, 44 ở $M^{+}$-30 và $V$-30; độ muộn lớn nhất của
các nhánh là 59,745 giây. Chúng không được miễn trừ bằng sổ hành vi: đây là
quy tắc thực thi riêng của phiên bản thí nghiệm. Vì vậy 0 quyết định miễn
trừ của $C$ không nói rằng mọi điều kiện đón nghiêm ngặt ban đầu đều thỏa.

Bảy trong tám dự đoán niêm phong của T38 đúng. Dự đoán sai yêu cầu
$M^{+}$-30 có không quá ba ô tối với khách vượt hạn chót, nhưng thực tế là
5/11. Việc dự đoán được niêm phong không loại nhu cầu báo các thất bại;
nó chỉ làm chúng có thể được kiểm và không sửa lại sau kết quả.

Kết luận của T38 có phạm vi hẹp: trên tập ô đã dựng, số khách của các cổng
30 giây gần nhau, cổng hành vi không cần miễn trừ, và khoảng dự phòng cố
định bao phần nhìn thấy trong hầu hết ô. Nó không sửa thất bại của H6,
không chứng minh độc lập giữa hiệu chỉnh và kiểm tra, không bảo đảm 90\%
cho quần thể, và không xác nhận một lời hứa 120 giây ngoài sổ hành vi.
